\paragraph{Observações.}
\begin{itemize}
  \item \textbf{O custo é quase todo do fator:} o transiente de 48~h leva 114--131~s (480 passos, 4,0 a 5,0 iterações
    de Newton por passo, sem nenhum corte de passo), e as 49 análises do fator levam 1700--2900~s (33--58~s por FS com
    o PC sobrecarregado). Os FS altos custam mais (o caminho de $F$ é mais longo): 53--58~s por FS nos casos com chuva
    fraca contra 33--38~s nos que rompem.
  \item \textbf{O refinamento adaptativo é indispensável e ainda não convergiu.} Na malha grossa o FS fica 0,05 a 0,20
    acima (f50 com 36~mm/h: 1,229 contra 1,034). Cada nível reduz a mudança mais ou menos à metade (no mesmo caso:
    $-0{,}110$, $-0{,}055$, $-0{,}020$, $-0{,}010$), e o último nível ainda muda o FS em 0,003 a 0,012. Se a mudança
    continuar caindo à metade, o FS convergido fica mais $\approx 0{,}01$ abaixo do relatado (Richardson). Para o
    estudo de P$_f$ isso é sistemático (todas as realizações caem juntas) e desloca a P$_f$: é preciso medir com
    \texttt{nref=5} (campanha do fermi, T4) e decidir entre refinar mais ou corrigir.
  \item \textbf{A malha do transiente é grossa} (catetos de 2,5~m) para uma frente de umedecimento de 1--2~m (f100);
    o FS do instante 0 já muda 0,005 com outra malha do transiente, porque o perfil inicial (hidrostático com teto de
    75~kPa) tem uma quina que a interpolação linear suaviza. Testada no fermi (T4: \texttt{tri3}, \texttt{tri4} e
    camadas de 0,25~m).
\end{itemize}
